% CPU
% CPU.tex

\documentclass[12pt,UTF8]{ctexbook}

% 设置纸张信息。
\input{../../../common/page}

% 设置字体，并解决显示难检字问题。
\xeCJKsetup{AutoFallBack=true}
\setCJKmainfont{SimSun}[BoldFont=SimHei, ItalicFont=KaiTi, FallBack=SimSun-ExtB]

% 目录 chapter 级别加点（.）。
\usepackage{titletoc}
\titlecontents{chapter}[0pt]{\vspace{3mm}\bf\addvspace{2pt}\filright}{\contentspush{\thecontentslabel\hspace{0.8em}}}{}{\titlerule*[8pt]{.}\contentspage}

% 设置 part 和 chapter 标题格式。
\ctexset{
	chapter/name={第,章},
	chapter/number={\chinese{chapter}}
}

% 图片相关设置。
\usepackage{graphicx}
\graphicspath{{Images/}}

% 设置署名格式。
\newenvironment{shuming}{\hfill\zihao{4}}

% 注脚每页重新编号，避免编号过大。
\usepackage[perpage]{footmisc}

\title{\heiti\zihao{0} CPU}
\author{佚名}
\date{}

\begin{document}

\maketitle
\tableofcontents

\frontmatter

\mainmatter

\chapter{第一章 CPU概述}

\section{CPU的定义与作用}
\section{CPU的发展历程}
\section{CPU的分类}

\chapter{第二章 CPU架构}

\section{冯·诺依曼架构}
\section{哈佛架构}
\section{现代CPU架构特点}

\chapter{第三章 CPU性能指标}

\section{主频与外频}
\section{缓存大小}
\section{核心数量}
\section{线程技术}
\section{功耗指标}

\chapter{第四章 CPU技术发展}

\section{早期CPU技术}
\section{多核技术发展}
\section{制程工艺进步}
\section{未来技术趋势}

\chapter{第五章 CPU应用领域}

\section{个人电脑}
\section{服务器与工作站}
\section{移动设备}
\section{嵌入式系统}
\section{超级计算}

\chapter{第六章 CPU选购指南}

\section{需求分析}
\section{预算考虑}
\section{品牌选择}
\section{性能对比}
\section{兼容性检查}

\chapter{第七章 CPU未来展望}

\section{量子计算}
\section{神经形态芯片}
\section{光子计算}
\section{新材料应用}

\chapter{第八章 CPU制造工艺}

\section{硅晶圆制备}
\section{光刻技术}
\section{蚀刻工艺}
\section{封装测试}

\chapter{第九章 CPU散热技术}

\section{风冷散热}
\section{水冷散热}
\section{液氮散热}
\section{热管技术}
\section{散热材料}

\chapter{第十章 CPU超频与优化}

\section{超频原理}
\section{超频风险}
\section{电压调节}
\section{稳定性测试}
\section{性能优化}

\chapter{第十一章 CPU故障诊断与维护}

\section{常见故障类型}
\section{诊断工具}
\section{维护方法}
\section{故障预防}
\section{维修建议}

\chapter{第十二章 CPU市场分析}

\section{市场份额}
\section{价格趋势}
\section{竞争格局}
\section{供应链分析}
\section{未来展望}

\chapter{第十三章 CPU与主板兼容性}

\section{接口标准}
\section{芯片组支持}
\section{BIOS设置}
\section{驱动程序}
\section{升级兼容性}

\chapter{第十四章 CPU缓存技术}

\section{L1缓存}
\section{L2缓存}
\section{L3缓存}
\section{缓存一致性}
\section{缓存优化}

\chapter{第十五章 CPU指令集架构}

\section{x86架构}
\section{ARM架构}
\section{RISC-V架构}
\section{MIPS架构}
\section{指令集对比}

\chapter{第十六章 CPU虚拟化技术}

\section{虚拟化原理}
\section{硬件虚拟化}
\section{软件虚拟化}
\section{容器技术}
\section{虚拟化性能}

\chapter{第十七章 CPU多核与并行计算}

\section{多核架构}
\section{并行编程}
\section{任务调度}
\section{负载均衡}
\section{性能优化}

\chapter{第十八章 CPU功耗与能效}

\section{功耗来源}
\section{能效指标}
\section{节能技术}
\section{动态频率调节}
\section{绿色计算}

\chapter{第十九章 CPU安全机制}

\section{硬件安全}
\section{加密技术}
\section{可信执行}
\section{漏洞防护}
\section{安全标准}

\chapter{第二十章 CPU测试与基准}

\section{性能测试}
\section{稳定性测试}
\section{兼容性测试}
\section{基准软件}
\section{测试方法论}

\chapter{第二十一章 CPU历史回顾}

\section{早期计算机}
\section{微处理器诞生}
\section{个人电脑革命}
\section{互联网时代}
\section{移动计算时代}

\chapter{第二十二章 CPU品牌对比}

\section{Intel}
\section{AMD}
\section{ARM}
\section{其他品牌}
\section{品牌选择建议}

\chapter{第二十三章 CPU与操作系统}

\section{Windows系统}
\section{Linux系统}
\section{macOS系统}
\section{移动系统}
\section{系统优化}

\chapter{第二十四章 CPU编程优化}

\section{编译器优化}
\section{代码优化}
\section{并行编程}
\section{缓存优化}
\section{性能分析}

\chapter{第二十五章 CPU未来技术趋势}

\section{3D芯片}
\section{异构计算}
\section{神经芯片}
\section{光子芯片}
\section{新材料}

\chapter{第二十六章 CPU故障排除}

\section{启动故障}
\section{运行故障}
\section{过热问题}
\section{兼容性问题}
\section{解决方案}

\chapter{第二十七章 CPU升级指南}

\section{升级评估}
\section{兼容性检查}
\section{升级步骤}
\section{性能提升}
\section{注意事项}

\chapter{第二十八章 CPU配件与接口}

\section{散热器}
\section{主板接口}
\section{内存接口}
\section{扩展插槽}
\section{连接标准}

\chapter{第二十九章 CPU工作原理}

\section{取指执行}
\section{指令解码}
\section{运算执行}
\section{结果写回}
\section{控制流程}

\chapter{第三十章 CPU寄存器与内存}

\section{通用寄存器}
\section{专用寄存器}
\section{内存层次}
\section{地址映射}
\section{数据传输}

\chapter{第三十一章 CPU流水线技术}

\section{流水线原理}
\section{流水线级数}
\section{流水线冒险}
\section{分支预测}
\section{乱序执行}

\chapter{第三十二章 CPU分支预测}

\section{静态预测}
\section{动态预测}
\section{预测算法}
\section{预测准确率}
\section{性能影响}

\chapter{第三十三章 CPU乱序执行}

\section{乱序原理}
\section{指令调度}
\section{执行单元}
\section{重排序缓冲}
\section{性能提升}

\chapter{第三十四章 CPU缓存一致性}

\section{一致性协议}
\section{MESI协议}
\section{缓存同步}
\section{伪共享}
\section{优化策略}

\chapter{第三十五章 CPU中断处理}

\section{中断类型}
\section{中断优先级}
\section{中断处理流程}
\section{中断延迟}
\section{实时性要求}

\chapter{第三十六章 CPU总线架构}

\section{前端总线}
\section{后端总线}
\section{系统总线}
\section{总线协议}
\section{总线性能}

\chapter{第三十七章 CPU时钟频率}

\section{时钟信号}
\section{频率倍频}
\section{时钟分配}
\section{时钟抖动}
\section{时钟同步}

\chapter{第三十八章 CPU核心电压}

\section{电压调节}
\section{动态电压}
\section{电压稳定性}
\section{功耗控制}
\section{电压监控}

\chapter{第三十九章 CPU封装技术}

\section{封装类型}
\section{引脚布局}
\section{散热设计}
\section{电气特性}
\section{封装趋势}

\chapter{第四十章 CPU制程工艺}

\section{制程节点}
\section{光刻技术}
\section{蚀刻工艺}
\section{材料选择}
\section{制程挑战}

\chapter{第四十一章 CPU硅晶圆制造}

\section{晶圆制备}
\section{晶圆切割}
\section{良品率}
\section{缺陷检测}
\section{成本控制}

\chapter{第四十二章 CPU缺陷与良品率}

\section{缺陷类型}
\section{缺陷检测}
\section{良品率计算}
\section{质量控制}
\section{改进措施}

\chapter{第四十三章 CPU热设计功耗}

\section{TDP概念}
\section{功耗测量}
\section{散热设计}
\section{能效优化}
\section{热管理}

\chapter{第四十四章 CPU集成显卡}

\section{核显架构}
\section{显存共享}
\section{图形性能}
\section{独显对比}
\section应用场景}

\chapter{第四十五章 CPU与AI计算}

\section{AI加速}
\section{神经网络}
\section{机器学习}
\section{深度学习}
\section{未来展望}

\chapter{第四十六章 CPU量子计算展望}

\section{量子比特}
\section{量子门}
\section{量子算法}
\section{量子优势}
\section{技术挑战}

\backmatter

《CPU》 佚名 著


\end{document}